% ARQUITECTURA_AUTORAL_PREEXISTENTE: semicuadrado bidireccional de N42.
% FORMALIZACION_NUEVA: grupoide de pares y ponderacion por estabilizadores.
\subsubsection{Incidencia marcada y normalización bidireccional}
\label{sec:ii09-normalizacion-bidireccional}

Los invariantes del lector conservan la naturaleza del objeto del que
proceden. La intersección marcada $B_K=H_e\cap P_\pi$ tiene cardinal
$b=4$; el entero $q=7$ es la coordenada del pivote $p_\varphi$ en la
numeración orientada. El orden de determinación del diseño de Steiner
es $p=5$. Los soportes hexádico y octádico tienen cardinales $h_6=6$
y $o_8=8$, y el alfabeto de orientación trítica tiene cardinal $t=3$.
Un cardinal de soporte y una coordenada de pivote se utilizan mediante
sus aplicaciones respectivas, aun cuando ambos se publiquen como enteros.

La normalización del autocruce pentagonal admite la siguiente
realización explícita. Sea $P$ el soporte de cinco posiciones del cierre
y sea $X=P\times P$ su espacio de pares ordenados. La inversión
bidireccional sobre este espacio es
\[
 \jmath:X\longrightarrow X,\qquad \jmath(x,y)=(y,x),
 \qquad\jmath^2=I.
\]
Se conserva la información de estabilizador de cada órbita. Para una
órbita $[z]$ se define su peso como
$1/|\operatorname{Stab}_{C_2}(z)|$. La suma de esos pesos es la
cardinalidad del grupoide de acción $X/\!/C_2$.

\begin{proposition}[Medida orbital del autocruce bidireccional]
\label{prop:ii09-semicuadrado-orbital}
Para un soporte finito de cardinal $p$, la inversión de pares ordenados
determina
\[
 |(P\times P)/\!/C_2|
 =\sum_{[z]\in(P\times P)/C_2}
      \frac1{|\operatorname{Stab}_{C_2}(z)|}
 =\frac{p^2}{2}.
\]
Para $p=5$, el autocruce tiene diez órbitas libres y cinco órbitas
fijas, y su medida orbital es $10+5/2=25/2$.
\end{proposition}
\begin{proof}
Cada par diagonal $(x,x)$ forma una órbita fija, cuyo estabilizador
tiene orden dos; hay $p$ de ellas. Los $p(p-1)$ pares restantes se
agrupan de dos en dos, con estabilizador trivial. Su contribución es
por tanto $p(p-1)/2$, y la de la diagonal es $p/2$. La suma es
$p^2/2$. Equivalentemente, la identidad órbita--estabilizador da
$1/|\operatorname{Stab}(z)|=|[z]|/2$; sumar sobre todas las órbitas
produce $|P\times P|/2$.
\end{proof}

La medida conserva explícitamente la contribución de los puntos fijos.
El conjunto ordinario de órbitas tiene cardinal $p(p+1)/2$, que vale
quince en el cierre pentagonal. La medida orbital y ese cardinal son
dos lecturas determinadas de la misma inversión. La normalización
bidireccional del lector utiliza la primera: las diez órbitas libres
pesan una unidad y las cinco diagonales pesan media unidad.

Esta realización precisa el factor $p^2/2$ empleado en el sistema
sectorial siguiente. La asignación de esa contribución al sector $23$,
su orientación negativa y las otras dos reglas de incidencia permanecen
especificadas por el sistema sectorial completo. La prueba establece la
normalización del autocruce y preserva diferenciada la selección de los
funcionales angulares.
